%=============================================================================!
% 12.1  PROBABILITY                                                           !
%=============================================================================!
\section{Probability}
\label{sec:sm-probability}

A liquid of a few hundred atoms passes through more arrangements in a
picosecond than could ever be listed, and what can be said about it is
what fraction of the time, or of many copies of it, has some property.
That is the language of probability. This section sets out the few rules
of it that the chapter uses, on a pair of dice first and then on a
quantity that varies continuously.

\begin{toolbox}[Probability]
A trial, such as a throw of two dice, has outcomes. Repeated many times,
the fraction of trials with a given outcome settles towards a number, its
probability $\mathcal{P}$, which lies between 0 and 1; the probabilities
of all the distinct outcomes add up to 1, since every trial has one of
them. When the outcomes are equally likely, a probability is a count: two
dice have $6\times6 = 36$ equally likely outcomes, of which $6 -
\lvert s - 7\rvert$ give the total $s$, for $s$ from 2 to 12: the first
die may show any face that leaves the second a face from 1 to 6, which is
$s - 1$ faces for $s \le 7$ and $13 - s$ for $s \ge 7$. So
$\mathcal{P}(s) = (6 - \lvert s - 7\rvert)/36$, largest at $s = 7$, with
$6/36$. Two outcomes of separate trials are independent if knowing one
does not change the probability of the other; then the probability of
both is the product of the two, as the 36 outcomes of two dice are the
$6\times6$ pairs, each with probability $\frac16\times\frac16$.
\end{toolbox}

\Cref{fig:sm-probability}(a) throws two dice by computer. After 100 throws
the fractions miss the probabilities by up to 0.051; after 10\,000 by up
to 0.0033. A fraction is only an estimate of a probability, and a better
one the more trials it rests on.

\begin{figure}[htb]
\centering
\includegraphics{ch12_ensembles/probability}
\caption{Frequencies settle on probabilities. (a) The fraction of throws
of two dice giving each total, after 100 throws (open) and 10\,000 throws
(filled), against the probabilities (bars). (b) The histogram of 1000
numbers drawn from the normal distribution, as a density, against the
Gaussian of mean 0 and standard deviation 1 (dashed).}
\label{fig:sm-probability}
\end{figure}

\subsection*{Averages and spreads}

\begin{toolbox}[Expectation and variance]
A quantity $x$ that takes the value $x_k$ with probability $\mathcal{P}_k$
has the expectation, or mean, $\ens{x} = \sum_kx_k\mathcal{P}_k$, the value
its average over many trials settles towards, and any function of it the
expectation $\ens{f(x)} = \sum_kf(x_k)\mathcal{P}_k$. The variance
measures the spread about the mean,
\begin{align*}
   \operatorname{var}(x) = \ens{(x - \ens{x})^2}
   &= \ens{x^2 - 2\ens{x}\,x + \ens{x}^2}\\
   &= \ens{x^2} - 2\ens{x}\ens{x} + \ens{x}^2 = \ens{x^2} - \ens{x}^2 ,
\end{align*}
expanding the square and then averaging term by term, which uses three
properties of $\sum_kf(x_k)\mathcal{P}_k$: the expectation of a sum is the
sum of the expectations, a constant factor such as $\ens{x}$ comes
outside, and a constant $c$ averages to $c\sum_k\mathcal{P}_k = c$. Its
square root is the standard deviation $\sigma_x$, in the units of
$x$ (written with the quantity as a subscript, to keep it apart from the
Lennard-Jones distance $\sigma$). For independent $x$ and $y$ the
probability of each pair is the product, so
\begin{equation*}
   \ens{xy} = \sum_{k,l}x_ky_l\,\mathcal{P}(x_k)\,\mathcal{P}(y_l)
   = \Bigl(\sum_kx_k\mathcal{P}(x_k)\Bigr)\Bigl(\sum_ly_l\mathcal{P}(y_l)\Bigr)
   = \ens{x}\ens{y} ,
\end{equation*}
and by the same expansion with $x + y$ in place of $x$,
\begin{align*}
   \operatorname{var}(x + y) &= \ens{x^2} + 2\ens{xy} + \ens{y^2}
   - \ens{x}^2 - 2\ens{x}\ens{y} - \ens{y}^2\\
   &= \operatorname{var}(x) + \operatorname{var}(y)
   + 2\bigl(\ens{xy} - \ens{x}\ens{y}\bigr)
   = \operatorname{var}(x) + \operatorname{var}(y) :
\end{align*}
the variances of independent quantities add.
\end{toolbox}

One die has the mean $(1 + 2 + \dotsb + 6)/6 = 3.5$ and the variance
$35/12$ (\cref{ex:sm-die}); two independent dice have twice each, 7 and
$35/6 = 5.8333$, and the 10\,000 throws above give 6.997 and 5.846.

\subsection*{Densities}

A velocity or a position takes a continuum of values, and the probability
of any exact value is zero. It has instead a probability density
$\mathcal{P}(x)$: the probability that $x$ lies between $a$ and $b$ is
$\int_a^b\mathcal{P}(x)\,\dd x$, so a narrow range of width $\Delta x$ at
$x$ holds the probability $\mathcal{P}(x)\,\Delta x$. Cutting all the
values into such ranges and letting them narrow, the sums above become
integrals: $\int\mathcal{P}\,\dd x = 1$ over all values, and $\ens{f(x)} =
\int f(x)\,\mathcal{P}(x)\,\dd x$. A histogram estimates a density: the fraction of
samples in a bin of width $\Delta x$, divided by $\Delta x$, approaches
$\mathcal{P}(x)$ as the samples grow many and the bins narrow. The most
important density of the book is the Gaussian, or normal, density,
\begin{equation}
   \mathcal{P}(x) = \frac{1}{\sigma_x\sqrt{2\pi}}\,
   \mathrm{e}^{-(x - \ens{x})^2/(2\sigma_x^2)} .
   \label{eq:sm-gauss}
\end{equation}
Its integral is 1, by the Gaussian integral $\int\mathrm{e}^{-au^2}\dd u =
\sqrt{\pi/a}$ of \cref{sec:md-ewald} with $u = x - \ens{x}$ and $a =
1/(2\sigma_x^2)$ (\cref{ex:md-normal}). It is symmetric about $\ens{x}$: the values
$\ens{x} + u$ and $\ens{x} - u$ have the same density, so their deviations
from $\ens{x}$ cancel in the average, and $\ens{x}$ is its mean. Its variance follows by integration by parts
(\cref{sec:la-euler}), with $u = x - \ens{x}$, $s = \sigma_x$, $f = u$
and $g' = u\,\mathrm{e}^{-u^2/(2s^2)}$, whose integral is $g =
-s^2\mathrm{e}^{-u^2/(2s^2)}$, as the chain rule confirms:
\begin{align*}
   \int u^2\mathrm{e}^{-u^2/(2s^2)}\dd u
   &= \Bigl[-s^2u\,\mathrm{e}^{-u^2/(2s^2)}\Bigr]_{-\infty}^{\infty}
   + s^2\int\mathrm{e}^{-u^2/(2s^2)}\dd u\\
   &= 0 + s^2\times s\sqrt{2\pi} ,
\end{align*}
and the bracket vanishes at both ends, where the exponential falls faster
than $u$ grows; dividing by $s\sqrt{2\pi}$, the variance is $\sigma_x^2$.
The fraction of the probability within $k$ standard deviations of the mean
follows with $w = u/(\sigma_x\sqrt2)$, so that $\dd u = \sigma_x\sqrt2\,\dd
w$ and the limits $\pm k\sigma_x$ become $\pm k/\sqrt2$:
\begin{equation*}
   \frac{1}{\sigma_x\sqrt{2\pi}}\int_{-k\sigma_x}^{k\sigma_x}
   \mathrm{e}^{-u^2/(2\sigma_x^2)}\,\dd u
   = \frac{1}{\sqrt\pi}\int_{-k/\sqrt2}^{k/\sqrt2}\mathrm{e}^{-w^2}\,\dd w
   = \operatorname{erf}(k/\sqrt2) ,
\end{equation*}
the last step since $\mathrm{e}^{-w^2}$ takes the same value at $w$ and
$-w$, with the error function of \cref{sec:md-ewald}: 68.27\%, 95.45\% and
99.73\% for $k$ = 1, 2 and 3.
The function \texttt{rng.normal} that drew the starting velocities of
\cref{sec:md-start} returns numbers with this density, mean 0 and
standard deviation 1; \cref{fig:sm-probability}(b) draws 1000 of them,
whose mean is 0.0029 and standard deviation 1.0115, and whose histogram
departs from \cref{eq:sm-gauss} by up to 0.042 per unit. Why velocities
should be Gaussian at all is the subject of \cref{sec:sm-maxwell}.

\begin{keyidea}
A probability is the long-run fraction of trials with an outcome; the
probabilities of independent outcomes multiply. A quantity has a mean
$\ens{x}$ and a variance $\ens{x^2} - \ens{x}^2$, and the variances of
independent quantities add. A continuous quantity has a density, which a
histogram estimates; the Gaussian density of \cref{eq:sm-gauss} has mean
$\ens{x}$ and variance $\sigma_x^2$.
\end{keyidea}

\notebook{12}{throw dice and draw normal numbers, and watch the histograms settle}

\begin{selfcheck}
What is the probability that two dice total at least 10?
\end{selfcheck}
